\documentclass[a4paper,11pt]{report}

% ── Encoding & fonts ──────────────────────────────────────────────────────
\usepackage[utf8]{inputenc}
\usepackage[T1]{fontenc}
\usepackage{lmodern}
\usepackage[margin=24mm]{geometry}

% ── Core packages ─────────────────────────────────────────────────────────
\usepackage{amsmath,amssymb}
\usepackage{enumitem}
\usepackage{booktabs}
\usepackage{array}
\usepackage{titlesec}
\usepackage{fancyhdr}
\usepackage{graphicx}
\usepackage{caption}
\usepackage[hidelinks]{hyperref}

% Shared brand style (palette + callout boxes). Relative path from textbook/.
\usepackage{../style/greatwall}

\graphicspath{{./}{../assets/}}

\hypersetup{colorlinks=true, linkcolor=accentdark, urlcolor=accentdark,
            citecolor=accentdark}

% ── Headers / footers ─────────────────────────────────────────────────────
\pagestyle{fancy}
\fancyhf{}
\renewcommand{\headrulewidth}{0pt}
\fancyhead[R]{\small\itshape\color{graytext}Great Wall — Pilot Course}
\fancyfoot[C]{\small\color{graytext}\thepage}

% ── Section styling ───────────────────────────────────────────────────────
\titleformat{\chapter}[display]
  {\bfseries\color{accent}}
  {\filright\large\color{graytext}\chaptertitlename\ \thechapter}
  {0.3em}
  {\Huge}
  [\vspace{-0.4em}{\color{accent}\rule{\textwidth}{1pt}}]
\titleformat{\section}{\large\bfseries\color{accentdark}}{\thesection}{0.6em}{}
\titleformat{\subsection}{\normalsize\bfseries\color{darktext}}{\thesubsection}{0.6em}{}

\newcommand{\probsoln}[3]{%
  \begin{description}[leftmargin=0pt,labelsep=0.5em,style=nextline]
    \item[\textcolor{attackcol}{Problem.}] #1
    \item[\textcolor{accent}{Solution.}] #2
    \item[\textcolor{modelcol}{Methodology.}] #3
  \end{description}}

\begin{document}

% ══════════════════════════════════════════════════════════════════════════
% TITLE PAGE
% ══════════════════════════════════════════════════════════════════════════
\begin{titlepage}
\centering
\vspace*{2cm}
{\fontsize{40}{46}\selectfont\bfseries\color{accent}Great Wall\par}
\vspace{6pt}
{\Large\color{darktext}A Practical, Hands-On Pilot Course\par}
\vspace{4pt}
{\large\itshape\color{graytext}Coercion-resistant self-custody through tacit knowledge\par}
\vspace{1.4cm}
{\color{accent}\rule{90mm}{1pt}\par}
\vspace{1.4cm}
\begin{minipage}{0.82\textwidth}\centering\small\color{graytext}
\itshape It's all in your head, in nobody else's; the attacker knows that, and
is nevertheless unable to rob it.
\end{minipage}
\vfill
{\small\color{graytext}Companion textbook to the five-session pilot.\\
Tracks chained protocol \texttt{0.3.0}. Draft for cohort~1 — expect bugs;
report them.\par}
\vspace{8pt}
{\small\color{graytext}Dual-licensed Apache-2.0 / MIT.\par}
\end{titlepage}

% ══════════════════════════════════════════════════════════════════════════
% PREFACE
% ══════════════════════════════════════════════════════════════════════════
\chapter*{How to use this book}
\addcontentsline{toc}{chapter}{How to use this book}
\markboth{}{}

This is the written companion to a live, hands-on course. It is deliberately
\textbf{succinct}: it carries the spine of each session so you can prepare
before and review after, but the real learning happens at the keyboard, in the
viewer, and in conversation with the cohort. Read the relevant chapter
\emph{before} its session.

\section*{The shape of every chapter}
Each chapter opens with the same three-part frame, and so does most of the
course:
\probsoln{What goes wrong — the specific failure we are trying to fix.}
{What Great Wall does about it.}
{How we will convince ourselves it actually works (usually: by doing it).}

\section*{The boxes}
\begin{keyidea}
The one sentence to carry out of a section. If you remember nothing else,
remember the key ideas.
\end{keyidea}
\begin{tryit}
A hands-on action — at the viewer, the CLI, or with a partner. The course lives
here.
\end{tryit}
\begin{deeper}
Optional depth for the curious. \textbf{Never required} to keep up — skip freely
and the main thread still makes sense.
\end{deeper}
\begin{pitfall}
A common misunderstanding, flagged so you can sidestep it.
\end{pitfall}
\noindent And the feedback marker:
\gwfeedback{Wherever you see this flag, we \emph{predict} this is a confusing
spot. If it confused you, say so loudly — that is exactly the signal the pilot
needs.}

\section*{What we assume of you}
You are comfortable with Bitcoin self-custody at the user level — you know what a
seed phrase is, why holding your own keys matters, and roughly what hashing and
signatures do. We do \emph{not} assume cryptography, fractal mathematics, Rust,
or information theory; every such idea is built from zero when first needed.

\begin{pitfall}[Safety — read this once, mean it forever]
Everything in this course uses \textbf{throwaway setups and test entropy only}.
Never put a real seed, or real funds, into a practice setup. Great Wall's own
rule holds in class too: the master secret is never shown, written, or shared —
not even with your instructor.
\end{pitfall}

\tableofcontents

% ══════════════════════════════════════════════════════════════════════════
\chapter{The Wrench in the Works}
% ══════════════════════════════════════════════════════════════════════════

\probsoln
{Self-custody hands you total control of your money — and therefore total
exposure to anyone who can physically reach you. A growing wave of
\wrench{}s ($\$5$-wrench attacks: ``give me your keys or else'') turns
``be your own bank'' into ``be your own vulnerable vault.''}
{A custody method whose secret \emph{cannot be handed over}, even under duress,
while keeping full individual self-custody and an openly published design.}
{Look at the trend and its price first (this chapter), then earn each property
over the next four sessions by building and using the thing.}

\section{The four-property test}
An adequate self-custody method should satisfy \emph{four} security properties at
once. Until now they have been mutually exclusive across deployed methods:
\term{Deviceless} (access tied only to your mind, not to a seizable object),
\term{Individual Custody} (no third party between you and your funds),
\term{Anti-Obscure} (security does not rely on the attacker not knowing the
trick), and \term{Coercion-Resistant} (force does not extract access within any
window an attacker can sustain). Score the market against them:

\begin{center}\small
\renewcommand{\arraystretch}{1.35}
\setlength{\tabcolsep}{4pt}
\begin{tabular}{@{}>{\raggedright\arraybackslash}p{0.235\textwidth} c c c c >{\raggedright\arraybackslash}p{0.30\textwidth}@{}}
\toprule
\textbf{Method} & \textbf{Devi.} & \textbf{Indiv.} & \textbf{Anti-obsc.} & \textbf{Coerc.} & \textbf{Fails on / cost} \\
\midrule
\textbf{TKBA} (Great Wall) & \gwcheck & \gwcheck & \gwcheck & \gwcheck &
only risk is \emph{forgetting} --- solved by the memory coach (Session~5) \\
Physical (vault, Shamir, multisig) & \gwcross & \gwcheck & \gwcheck & \gwcheck &
astronomical cost + geographic binding: ``glorified gold'' \\
Shared / delegated (Coinbase, Casa) & \gwcheck & \gwcross & \gwcheck & \gwcheck &
negates self-custody --- ``not your keys\ldots'' \\
Obscurity (decoy, denial, redirect) & \gwcheck & \gwcheck & \gwcross & \gwskull &
\emph{backfires}: turns robbery into torture / homicide \\
Vanilla (Ledger, Trezor, Electrum) & \gwcheck & \gwcheck & \gwcheck & \gwcross &
wide open to the wrench \\
\bottomrule
\end{tabular}
\end{center}
\begin{center}\footnotesize\color{graytext}
\gwcheck~has the property \quad \gwcross~lacks it \quad
\gwskull~\emph{negative} --- actively worse than merely lacking it (see below)
\end{center}

\begin{keyidea}
The four properties are a \term{covering set}, not a wish-list: drop any one and
the failure mode is concrete and named. Only \GW{} (TKBA) holds all four. The
recurring hole is \term{coercion-resistance} --- conventional tools guard against
\emph{remote} theft but do little against a person with a wrench.
\end{keyidea}

\noindent Read the obscurity row carefully: its coercion cell is a \emph{skull},
not a cross. Obscurity does not merely \emph{fail} to resist coercion --- it
\emph{inverts}. Once ``deny you own Bitcoin'' and ``hand over a decoy'' are the
community's standard advice, every signal a victim can emit has been
pre-discredited by the defense itself (we return to this next section). That is
worse than a vanilla wallet, whose owner can at least comply and survive.

\section{Obscurity is a cost, not a defense}
The community's consensus advice is obscurity: \emph{``No one needs to know you
own Bitcoin.''} It is offered as security; it is really a recurring tax. Two
problems:

\begin{enumerate}[leftmargin=1.4em,nosep]
\item \textbf{It contradicts adoption.} If safety requires being a needle in a
haystack, then everyone cannot be safe: \emph{if everyone is a needle, there is
no hay}. Obscurity and mass adoption are mutually exclusive — ``maximalism vs.\
obscurity: choose one.''
\item \textbf{It violates Kerckhoffs (1883).} A system's security must not depend
on the secrecy of its mechanism. Bitcoin is open by design; an educated attacker
shares the defender's knowledge, so obscurity-based tricks (decoy wallets,
redirect scripts) \emph{backfire} — the attacker simply keeps coercing past the
first ``compliant'' transfer, because they cannot tell a decoy from the real
stash.
\end{enumerate}

\subsection{The security tax, priced}
We can put a number on this. Insurance markets price physical-vulnerability risk
every day for portable luxury goods (jewelry, watches). Cross-referencing the
documented wrench-attack rate against jewelry premiums, the
\emph{insurance-equivalent} annual security cost for individual Bitcoin whales
has, at the 2026 pace, entered the \textbf{gold-storage range} ($\sim$0.5–1.5\%/yr)
— and, because attacks are growing faster than the whale population, the
\emph{forward-looking} discount over a multi-year horizon is several times the
current annual cost. The full derivation is in
\texttt{justification-and-economics/JUSTIFICATION}; for our purposes the
takeaway is blunt:

\begin{keyidea}
The wrench attack is not a tail-risk anecdote. It is a measurable, compounding
holding cost that the market currently prices at zero — and obscurity, the
status-quo answer, is part of the cost, not the cure.
\end{keyidea}

\section{Rhetoric can't save you: two cases and a trilemma}
\label{sec:trilemma}
Two documented attacks bracket the failure of \emph{everything except}
coercion-resistance. Put them side by side.

\subsection{Force beats a fortress: David Balland}
David Balland co-founded Ledger — a hardware-wallet company. If wealth,
technical mastery, and security awareness could deter a \wrench{}, he was as
well-defended as an individual gets. In January 2025 he was nonetheless
kidnapped and held for roughly two days, and mutilated, while his captors
extorted a cryptocurrency ransom; he was freed only by a police
operation.\footnote{One of the most protracted incidents in Jameson Lopp's
compilation, \emph{Known Physical Bitcoin Attacks},
\url{https://github.com/jlopp/physical-bitcoin-attacks}.} The lesson is not
``get more bodyguards.'' It is that \emph{abundant physical defense does not
prevent the attempt or the harm}: coercion-resistance has to live at the
\emph{protocol} layer, where there is nothing to hand over, not at the
operational-security layer, which force simply overruns.

\subsection{Denial gets you tortured: RATIRL}
The Swedish \emph{League of Legends} streamer RATIRL was bound and tortured for
hours by armed intruders demanding cryptocurrency he says he did not
have; unable to credit the denial, they persisted long past any plausible payoff,
and left with his Pokémon-card collection.\footnote{J.\ Lopp, on X: ``\ldots
armed intruders broke into his home, tortured him for hours, and demanded
cryptocurrency he didn't have. The robbers eventually left with only his
Pok\'emon cards\ldots'' \url{https://x.com/lopp/status/2058917155530293600}.} By
the victim's own account the ordeal was \emph{prolonged} — and that detail is the
tell. Prolonged coercion following a denial all but confirms the attackers were
\emph{predisposed to disbelieve} it, truthful or not. Denial did not save him;
it very likely \emph{extended} the attack.

\subsection{The rhetoric trilemma}
Why did denial fail? Because under coercion everything a victim can \emph{do} is
a \emph{speech act} whose credibility the \emph{attacker} adjudicates — and there
are only three such moves:

\begin{enumerate}[leftmargin=1.6em,nosep]
\item \term{Comply} — hand over the real key. Outcome: robbed. And if the
attacker suspects there is more, compliance does not even end the encounter.
\item \term{Deny} — ``I have nothing.'' Outcome: in a culture where ``deny you
own Bitcoin'' is standard advice, denial is pre-discredited as \emph{the
denial-script} → disbelief → escalation (RATIRL).
\item \term{Decoy} — surrender a small or fake wallet. Outcome: presumed to be
\emph{the decoy} → ``now the real one'' → escalation.
\end{enumerate}

\begin{keyidea}[The trilemma]
All three moves are \emph{unverifiable claims}, and obscurity culture has already
discredited the two non-complying ones. So the honest have-nothing victim and the
stonewalling holder look \emph{identical} to the attacker — whose rational
continuation of an interrogation that opened with a denial is \emph{more}
violence, not less. Rhetoric offers no exit.
\end{keyidea}

\noindent \GW{} dissolves the trilemma by changing the \emph{kind} of signal the
victim emits — from \emph{rhetoric} (a claim to be believed) to a
\emph{verifiable physical constraint}. ``I cannot produce it within your window''
is not something the attacker must take on faith: it is a property of a
\emph{public} protocol they can check for themselves (this is what
\term{Anti-Obscurity} buys — the one honest signal becomes credible \emph{because}
the attacker knows the mechanism). Balland's limb says the barrier must sit at the
protocol layer; RATIRL's limb says the victim's signal must be verifiable, not
rhetorical. Great Wall satisfies both at once.

\gwfeedback{The trilemma is subtle: students often first think ``just deny.''
RATIRL is the counter-example that makes it click — watch whether the
``pre-discredited signal'' idea lands or needs a second pass.}

\section{What we will build instead: the Four Properties}
\GW{} claims four properties \emph{at once}:

\begin{enumerate}[leftmargin=1.6em,nosep]
\item \term{Knowledge-Based Authentication.} The secret lives entirely in your
memory — no device, vault, or place required.
\item \term{Individual Custody.} You depend on no one else; full self-custody is
intact.
\item \term{Non-Obscurity.} The method is fully public; it does not fail when the
attacker learns it, and does not rely on hiding that a stash exists.
\item \term{Coercion-Resistance.} Threatening violence does not get the attacker
the secret.
\end{enumerate}

\begin{keyidea}
In one sentence: \emph{it's all in your head (1), in nobody else's (2), the
attacker is aware of that (3), and is nevertheless unable to rob it (4).}
\end{keyidea}

\noindent This sounds like having your cake and eating it too. The next four
sessions explain why it is nonetheless possible — the secret is \emph{tacit}
(Session~2), encoded as a place on an indescribable fractal (Session~3), behind a
wall of unavoidable computation (Session~4), made livable by training and
inheritance (Session~5).

\gwfeedback{Does the economic framing read as evidence, or as a sales pitch?
We need to know — this is the hook the whole course hangs on.}

\begin{tryit}[Homework into Session 2]
Install \texttt{great-wall-core}, build the Rust engine, and launch
\texttt{python3 viewer.py}. Then just \emph{play}: pan, zoom, change color
schemes. No goal. Get a feel for how strange and detailed the Burning Ship is.
Post a screenshot in \texttt{\#show-your-fractal}.
\end{tryit}

\paragraph{Reflect.} Write one paragraph on \emph{your} threat profile: who could
plausibly coerce you, for how long, and which conventional defense you currently
rely on. We build on this next time.

% ══════════════════════════════════════════════════════════════════════════
\chapter{Secrets You Cannot Tell}
% ══════════════════════════════════════════════════════════════════════════

\probsoln
{A secret you can \emph{say} is a secret you can be \emph{made} to say. Under a
long enough coercion, PINs, passphrases, and seed words all come out — and
``hidden extra'' tricks just get the victim hurt for plausibly holding more.}
{Encode the secret as \term{tacit knowledge} — something you can deploy
(recognize, point at) but genuinely \emph{cannot articulate} — and gate its use
behind a computation no one can rush.}
{Feel it before you define it: decode a point by hand, then try to describe its
location precisely to a partner. You will fail. That failure is the whole idea.}

\section{Explicit vs.\ tacit knowledge}
\term{Explicit knowledge} is anything you can state: a number, a word, a
sentence. It is transmissible — which is exactly why it is coercible. \term{Tacit
knowledge} is knowing-\emph{how} that resists being put into words: recognizing a
friend's face, riding a bicycle, the ``feel'' of a familiar route. You can use it
flawlessly and still be unable to dictate it to a stranger.

\begin{keyidea}
\textbf{Recognition is not recall, and recall is not verbalizable.} You can pick
your fractal point out of the landscape (recognition) without being able to read
out its coordinates (verbalization). Coercion can only extract what can be
verbalized.
\end{keyidea}

\section{Why this defeats the wrench}
A \wrench{} works by raising the cost of \emph{refusing} until you comply.
It assumes the secret is something you \emph{could} hand over if sufficiently
motivated. Tacit knowledge breaks that assumption: there is no sentence you can
be tortured into saying that conveys ``the point is \emph{there}.'' The only way
to extract it is to sit you in front of the fractal and have you point — which,
as we will see in Session~4, the attacker cannot make happen quickly, because the
fractal you would point at \emph{does not exist yet} until a long computation
finishes.

\begin{pitfall}
Tacit $\neq$ secret-in-the-usual-sense. The attacker may know \emph{everything}
about the method and that you have a stash (non-obscurity!). What they cannot get
is the non-transmissible \emph{recognition} in your head.
\end{pitfall}

\section{TKBA: the theory of the whole system}
\GW{} is an implementation of \term{Tacit Knowledge-Based Authentication
(TKBA)}. The four properties are logically coupled: a secret held only in your
head (1), by nobody else (2), known by the attacker to be there (3), can resist
coercion (4) \emph{only if} it is non-transmissible — i.e.\ tacit. TKBA is the
load-bearing premise, not an implementation detail.

TKBA needs a second ingredient: the \emph{interface} that deploys the secret must
itself be gated by an \term{inescapably lengthy computation}. Otherwise an
attacker just stands over you repeating ``try again'' until recognition leaks.
That gate is Argon2 (Session~4).

\begin{deeper}[Why a design choice can break TKBA]
Any feature that converts tacit recall into an \emph{explicit, memorizable fact}
weakens property~4. This is why the UX hides the bisection-area rectangles behind
a debug flag: showing them in normal use would teach you verbalizable facts about
your own encoding — facts an attacker could then coerce. Security here is a
property of \emph{what you are allowed to know about your own secret}.
\end{deeper}

\gwfeedback{Did ``tacit'' actually click, or stay abstract? The pairing exercise
is the moment of truth — watch whether students can articulate why their
description failed.}

\begin{tryit}[In session — the pairing exercise]
With a partner on Discord: each of you decodes one guided point in the viewer.
Now \emph{describe your point's location to your partner in words} well enough for
them to click it — no screen-share. Notice how quickly you run out of language,
and how ``up and a bit left, near that flamey bit'' fails to pin a 32-bit
location.
\end{tryit}

% ══════════════════════════════════════════════════════════════════════════
\chapter{Drawing Inside the Labyrinth}
% ══════════════════════════════════════════════════════════════════════════

\probsoln
{``Encode a secret as a point on a fractal'' sounds like hand-waving. How is the
mapping \emph{exact}, \emph{lossless}, and \emph{reversible} — and why a fractal
at all?}
{A \term{bisection bijection}: repeatedly split a rectangle, let the fractal's
own structure choose where, and read each input bit as a directional choice. Same
structure on both sides $\Rightarrow$ encode and decode walk the identical tree
$\Rightarrow$ perfect round-trip.}
{Do a tiny one by hand (2–3 bits), then scale up in the viewer and watch your
exact BIP39 words come back out.}

\section{First, some questions about bits and words}
The whole chapter — and the sizes table in Chapter~1 — leans on vocabulary we
have used without defining: \emph{bit}, \emph{128}, \emph{11 bits per word},
\emph{checksum}, and \emph{32 per point}. Settle them now, in order.

\subsection{What exactly is a ``bit of entropy''?}
Not a bit of \emph{storage} — a bit of \emph{unpredictability}. One bit of
entropy is one genuinely unpredictable coin-flip: it \emph{doubles} the number of
equally-likely secrets an attacker must sift. So $n$ bits of entropy means $2^n$
equally-likely possibilities, and on average $\approx 2^{n-1}$ guesses to hit
yours. The word \emph{unpredictable} is load-bearing: a bit that is always $1$
stores data but carries \emph{zero} entropy. Entropy counts only the part an
attacker cannot anticipate.

\subsection{Why do I need 128 of it?}
Because brute-force cost grows as $2^n$, and $2^{128}\approx 3.4\times10^{38}$ is
the line between ``someday, maybe'' and ``not before the heat death of the
universe.'' At 64 bits a well-funded adversary can grind through the whole space;
at 128 no one can, with any conceivable hardware, ever; past 128 you are spending
memory and effort for security you already had. It is also exactly BIP39's
default — a 12-word seed carries 128 bits of entropy (the 12 words actually spell
out 132 bits: 128 of entropy plus a 4-bit checksum — the next two questions).

\subsection{How does a BIP39 word encode 11 bits?}
The BIP39 wordlist is exactly \textbf{2048 words}, and $2048 = 2^{11}$. So a word
is nothing more than an \emph{index} into that list — a number from 0 to 2047 —
and any number in that range is written in \textbf{11 bits}. Twelve words
therefore carry $12\times 11 = \textbf{132}$ bits. (The words are a
human-friendly \emph{spelling} of the bits, chosen so the first four letters are
unique and easy to transcribe — nothing more mysterious than that.)

\subsection{What are checksum bits?}
Of those 132 bits, only \textbf{128 are entropy}; the last \textbf{4 are a
checksum}. The checksum is not random — it is \emph{derived} from the entropy: it
is the first $128/32 = 4$ bits of $\mathrm{SHA\text{-}256}$ of the entropy. Being
derived, it carries \emph{zero} entropy (it tells an attacker nothing). Its job is
\term{error detection}: when you type a seed, a wallet re-derives the checksum and
rejects the phrase if it disagrees, so a single mistyped or swapped word is
caught. A randomly chosen set of 12 valid words passes only $1/2^{4} = 1/16$ of
the time — that is the check working. (Bigger seeds get proportionally more:
24 words $=$ 256 entropy $+$ 8 checksum $= 264 = 24\times 11$.)

\subsection{How does one Great Wall point yield 32 of it?}
A point is pinned down by \textbf{32 binary choices} — the 32 left/right (or
up/down) bisection steps that locate its leaf (next section). If you picked that
leaf \emph{uniformly at random} from the $2^{32}$ possible leaves, then ``which
leaf?'' is worth exactly 32 bits. One point $=$ 32 bits — so 128 bits needs
\textbf{four} points, which is precisely why a 12-word setup has four point stages
(Chapter~1's table): $4\times32 = 128$. Great Wall encodes only those \emph{128
entropy} bits; the 4 checksum bits are \emph{recomputed} on decode, not stored on
the fractal (which is why it is four points, not ``$132/32$'').

\begin{keyidea}
Entropy is guessing-difficulty in bits ($n$ bits $\Rightarrow 2^n$ tries). 128
bits is the uncrackable bar. A BIP39 word is an index into a $2^{11}$-word list
$=$ 11 bits, so 12 words $=$ 132 bits $=$ 128 entropy $+$ 4 (derived, zero-entropy)
checksum. One Great Wall point is 32 uniformly-random bisection choices $=$ 32
bits, so four points carry the 128 entropy bits (the checksum is recomputed, not
stored). The \term{haystack} (which fractal you are on) is \emph{not} counted
here — that is the coercion cost (Session~4), not seed entropy. Entropy lives in
the \emph{needle}; the wall lives in the \emph{haystack}.
\end{keyidea}

\gwfeedback{The needle-vs-haystack split confuses people: they expect ``the
fractal'' to add entropy. It doesn't — it adds \emph{cost}. Watch for this.}

\section{The big picture: a shared, unstored codebook}
Both the encoder and the decoder can regenerate the \emph{same} fractal from the
same parameters. The fractal's intricate structure acts as a codebook neither
side has to store — they each reconstruct it on demand. Encoding writes a path
into that structure; decoding reads it back.

\begin{keyidea}
\textbf{Haystack $\gg$ needle.} Specifying \emph{which} fractal (the haystack)
is made far more expensive than specifying \emph{where} in it your point sits
(the needle, $\approx$ 32 bits). That asymmetry is the security engine: an
attacker who tries to \emph{describe} a fractal instead of \emph{deriving} it
faces the haystack cost, per stage, again and again.
\end{keyidea}

\section{Bisection encoding, step by step}
Picture the fractal inside a rectangle. To encode a bit string:

\begin{enumerate}[leftmargin=1.5em,nosep]
\item \textbf{Find structure.} Discover ``islands'' — connected regions sharing
an escape count — inside the current rectangle.
\item \textbf{Pick a split line.} Take a (weighted) median of the islands along
the rectangle's longer axis. The \emph{fractal} chooses the split, not a fixed
grid.
\item \textbf{Read one bit as a direction.} On a vertical split, $0=$ left,
$1=$ right; on a horizontal split, $0=$ up, $1=$ down. Select that child
rectangle.
\item \textbf{Contract if needed} (to stop one side ballooning), then
\textbf{recurse} on the chosen child with the next bit.
\end{enumerate}

\noindent After all bits, the \emph{center of the final rectangle} is your
encoded point. Decoding reproduces the very same splits and, at each level, asks
``which side is the point on?'' to read the bit back.

\begin{center}
\begin{tikzpicture}[scale=1.0, every node/.style={font=\footnotesize}]
  % level 0 rectangle
  \draw[accentdark, thick] (0,0) rectangle (6,3.2);
  % first split (vertical) at x=3.6
  \draw[attackcol, thick, dashed] (3.6,0) -- (3.6,3.2);
  \node[accentdark] at (3.0,3.45) {split 1 (median)};
  \node at (1.8,2.7) {\textbf{0}=left};
  \node at (4.8,2.7) {\textbf{1}=right};
  % chosen left child, second split (horizontal) at y=1.4
  \draw[accent, thick] (0,0) rectangle (3.6,3.2);
  \draw[attackcol, thick, dashed] (0,1.4) -- (3.6,1.4);
  \node at (1.0,2.1) {\textbf{0}=up};
  \node at (1.0,0.6) {\textbf{1}=down};
  % final leaf + point
  \draw[goldcol, very thick] (0,0) rectangle (3.6,1.4);
  \fill[btccol] (1.7,0.7) circle (2.2pt);
  \node[btccol] at (2.55,0.7) {point};
\end{tikzpicture}
\captionof{figure}{Two bisection levels. The fractal's island structure (not
shown) picks each dashed split line; each input bit selects a child; the leaf's
center is the encoded point. Decoding walks the identical tree and reads the
bits back from which side the point lies on.}
\end{center}

\section{Why it is genuinely a bijection}
Three guarantees, together:
\begin{itemize}[leftmargin=1.4em,nosep]
\item \textbf{Determinism.} The island discovery is driven by a pseudo-random
generator seeded from the \emph{bisection path so far} — so encode and decode,
walking the same path, discover identical islands and pick identical splits.
\item \textbf{No floating point.} Every bijection-critical step uses integer /
fixed-point arithmetic (\term{I4F60}: 60 fractional bits), so the result is
\emph{bit-identical} on any machine. Floats, with their platform-dependent
rounding, would silently break the round-trip.
\item \textbf{Clean partition.} Half-open intervals $[\min,\max)$ split a parent
exactly into two non-overlapping children, so every point lands in exactly one
leaf, and distinct bit strings land in distinct leaves.
\end{itemize}

\begin{pitfall}
``A fractal is fuzzy/chaotic, so this must be approximate.'' No. The chaos is in
the \emph{shape}; the \emph{arithmetic} is exact integer math. Same inputs,
identical outputs, forever. That exactness is why your seed survives the
round-trip.
\end{pitfall}

\gwfeedback{Where does the bijection lose people — the splitting, the
determinism (``why no floats''), or the ``why a fractal at all''? Pin which one.}

\begin{deeper}[Precision budget]
Each bisection level consumes roughly one bit of positional precision. With 60
fractional bits the encoder reliably places \textbf{32 bits per point}. That is
why the protocol uses exactly one 32-bit point per stage; bigger seeds use
\emph{more stages}, not more bits per point (see Session~4 and the size table in
the appendix).
\end{deeper}

\begin{tryit}[In session]
Encode a \textbf{3-word (32-bit) mini-seed} in the viewer. Turn on the area
visualization and step through the bisection (\texttt{V}, then \texttt{<}/\texttt{>}).
Then decode the point back and confirm you get the identical three words. The
skeptic's reassurance: your exact words return.
\end{tryit}

% ══════════════════════════════════════════════════════════════════════════
\chapter{Building the Wall}
% ══════════════════════════════════════════════════════════════════════════

\probsoln
{If the next fractal could be rendered \emph{instantly}, coercion would simply
wait for it: ``point at it now.'' What makes the secret unreachable \emph{in a
hurry} — even by a fully cooperative you?}
{A \term{memory-hard chain}. Each stage's fractal is derived by a slow, RAM-bound
Argon2 computation over stage-0 text plus all preceding points. It cannot be
traded for speed, and it is chained — stage $k{+}1$'s fractal does not exist until
stage $k$'s point is fixed. The wall-clock cost compounds.}
{Run a (small) calibration live; watch RAM fill and seconds pass; reason through
why no shortcut exists.}

\section{The wall in one picture}
\begin{center}
\begin{tikzpicture}[node distance=8mm, font=\footnotesize,
  box/.style={draw=accentdark, rounded corners=2pt, fill=accentlight,
              minimum height=8mm, inner sep=4pt, align=center},
  pt/.style={draw=btccol, thick, rounded corners=2pt, fill=white,
             minimum height=8mm, inner sep=4pt, align=center},
  arg/.style={draw=attackcol, thick, fill=white, rounded corners=2pt,
              inner sep=4pt, align=center},
  ar/.style={-latex, thick, accentdark}]
  \node[box] (s0) {stage 0\\\scriptsize text (salt/pepper)};
  \node[arg, right=of s0] (a1) {Argon2$^{N}$\\\scriptsize $\to$ SHA-256 $\to (o,p,q)$};
  \node[pt, right=of a1] (p1) {stage 1\\\scriptsize point};
  \node[arg, right=of p1] (a2) {Argon2$^{N}$};
  \node[pt, right=of a2] (p2) {stage 2\\\scriptsize point};
  \node[right=4mm of p2] (dots) {$\cdots$};
  \draw[ar] (s0) -- (a1);
  \draw[ar] (a1) -- (p1);
  \draw[ar] (p1) -- (a2);
  \draw[ar] (a2) -- (p2);
  \draw[ar] (p2) -- (dots);
  % feedback chain arrow: prior points feed the next Argon2
  \draw[ar, dashed] (p1) to[bend left=28] (a2);
\end{tikzpicture}
\captionof{figure}{The chain. Stage-0 text seeds it; each Argon2 pass consumes
stage-0 text plus \emph{all} preceding points to derive the next fractal's
parameters. There is no public ``canonical'' fractal — every point-bearing
surface is private and chain-derived.}
\end{center}

\section{What ``memory-hard'' buys}
\term{Argon2} is a key-stretching function that deliberately needs a lot of RAM
and runs as a strictly sequential chain of memory-dependent steps. Two
consequences:
\begin{itemize}[leftmargin=1.4em,nosep]
\item \textbf{You cannot trade memory for speed.} An attacker with a big GPU/ASIC
gains little: the work is gated on memory bandwidth and sequential dependency,
where commodity CPUs are already near-optimal.
\item \textbf{You cannot parallelize within a stage}, and you cannot start the
\emph{next} stage early — the chain forbids it. So total honest cost is
``one heavy Argon2 run per stage,'' and an attacker faces \emph{at least} that
much wall-clock to reach the last point.
\end{itemize}

\begin{keyidea}
Coercion fails twice over: the locations cannot be \emph{verbalized} (tacit), and
the later fractals cannot be \emph{materialized} faster than the full chain runs.
Robbing the whole secret needs custody of the victim for at least the total
Argon2 delay — and real wrench attacks do not last that long.
\end{keyidea}

\section{The design choices, in plain words}
You do not need the cryptographic proofs, but you should be able to say
\emph{why}:
\begin{description}[leftmargin=0pt,style=nextline,nosep]
\item[Stage 0 = salt \emph{or} pepper.] A short \texttt{[A-Z0-9-]} label.
As a \term{salt} (\texttt{RETIREMENT}, \texttt{MAIN-STASH}) it forks the
\emph{entire} chain — one memorized seed yields independent wallets per label,
and the memory of one does not transfer to another. As a \term{pepper} it lets
you stack this setup on top of a prior secret (a password manager entry, a Shamir
share, an earlier \GW{} setup).
\item[$p=1$ (no parallel lanes).] Here the delay \emph{is} the product, not a
budget to fill. One sequential lane maximizes the memory gate and keeps the
attacker/defender speed ratio near~1.
\item[$t=2$ (two passes).] The minimum that forces the whole memory array to stay
resident at once (the second pass re-reads it data-dependently), defeating the
classic time--memory trade-off.
\item[Argon2d (data-dependent).] Side-channel resistance is irrelevant on your
own device; data-dependent addressing gives the strongest memory-hardness, which
is what we want.
\item[Remember $N$, not ``hours.''] The durable parameter is the iteration count
$N$. ``Three days'' is only a \emph{perishable label} on $N$ — faster hardware
later changes how long a given $N$ takes, never the result. Recovery reproduces
the digest from $N$.
\end{description}

\gwfeedback{$p=1$ is the classic trip-wire: people trained on password hashing
``know'' more parallelism is better. The inversion (delay is the product) is the
crux — watch for pushback here.}

\begin{pitfall}[Checkpoints are a loaded gun]
A ``save intermediate'' option can store the chain's intermediate digests to
resume later — but that file is \emph{exactly} the artifact that lets an attacker
skip the time cost. It is \textbf{off by default}, and if you turn it on, it is
\emph{your} responsibility to securely delete it (\texttt{shred}/\texttt{srm})
once done. Sibling tools (TLP gating) automate this; the core exposes the raw
trade-off on purpose.
\end{pitfall}

\begin{tryit}[In session]
Pick a short target duration and let the app \emph{calibrate} $N$ on your device.
Watch memory usage climb. Then try a nearby iteration count using checkpoints to
feel how ``approximate memory of $N$'' is recoverable — and finish by securely
deleting the checkpoint file.
\end{tryit}

% ══════════════════════════════════════════════════════════════════════════
\chapter{Living With It}
% ══════════════════════════════════════════════════════════════════════════

\probsoln
{``Stake your retirement on remembering fractals'' invites three fair fears:
\emph{I'll forget}; \emph{the wall is too slow for me too}; and \emph{what
happens when I die or can't recall?}}
{Three companion systems: \term{Train} (measured memory consolidation with a
graduation gate), \term{Accelerate} (outsourceable time-lock puzzles without
losing custody), and \term{Inherit} (a dead-man's switch with cascading
fallbacks).}
{Do a complete setup once, then walk each ongoing obligation and decide, out
loud, whether it convinces you.}

\section{Train: don't trust memory you haven't measured}
\GW{} does not ask you to \emph{hope} you remember. The
\term{Celestial Peace: Never Forget} trainer drills your recall with spaced
repetition (SM-2, Anki-style). Crucially there is a \term{graduation gate}: your
stash is \textbf{not considered secured} until the trainer \emph{measures} that
your memory has consolidated. Until then you are bootstrapping.

\begin{keyidea}
Honesty is a feature. The hard question — ``what if I forget?'' — is answered by
\emph{measuring} consolidation before declaring the setup safe, not by
reassurance.
\end{keyidea}

\begin{deeper}[The bootstrap paradox, and the provisional key]
To \emph{secure} a stash you must already hold its entropy tacitly; to
\emph{memorize} it you must train against the setup over weeks. Neither precedes
the other, so a \term{provisional key} is held externally during the
consolidation window and \emph{destroyed at graduation}, after which the setup
self-secures. It is a transient crutch, not a backup.
\end{deeper}

\section{Accelerate: buy back time without giving up custody}
Set your Argon2 delay conservatively (say $2\times$), then tune the
\emph{effective} per-session delay with a \term{time-lock puzzle (TLP)} — an
RSW repeated-squaring puzzle whose setup is instant but whose solving takes a
chosen wall-clock time. You can \emph{outsource} the squaring to the
\term{jade-clock} Lightning marketplace: the solver receives only the puzzle, never
your vault, learns nothing about what it unlocks, and you pay over Lightning for
anonymity. Custody never leaves your device.

\section{Inherit: a phoenix, not a lawyer}
\term{phoenix-scroll} is a dead-man's switch. You periodically ``rotate'' a
private Lightning channel whose heir-branch key is split into two shares: one the
heir derives from \emph{their own} \GW{} recall, one you hand off each epoch
wrapped in a TLP. While you keep rotating, any heir's in-progress solve is
invalidated — they cannot front-run a living you. When rotation \emph{stops}
(death, incapacity), the latest puzzle eventually solves and the heir claims.
Opaque taproot fallbacks let inheritance \emph{cascade} (your heir's heirs)
without you ever knowing their estate plans.

\begin{pitfall}[The one genuinely terminal failure]
Without inheritance configured, \emph{total} loss of your tacit recall (with no
heir set up) is unrecoverable — by anyone, for the same reason coercion fails.
That is the price of true self-custody with no backdoor. Setting up inheritance is
the \emph{only} post-loss recovery path.
\end{pitfall}

\section{The feel of it: an instrument, not an app}
Two UX north-stars frame the whole experience: \emph{sober, but game-like}
(tactile, quietly rewarding, never gimmicky) and \emph{an instrument, not an app}
(you operate and master it; it does not think for you). The deeper bet: spatial
memory — the \emph{method of loci}, memory palaces — is the most robust trainable
human memory, and the fractal-as-territory is literally that. The faculty is
dormant, not absent.

\gwfeedback{After doing one setup, does ``trust my own memory'' feel plausible —
and what is the single biggest remaining doubt? This is the most important data
point in the whole pilot.}

\begin{tryit}[In session — the capstone]
Run a \textbf{complete setup} end-to-end (use a small/sub-standard size for time).
Configure a mock training schedule. On paper, sketch a one-level inheritance will
(one heir). Then take the skeptic's seat: name one assumption you are \emph{not}
yet willing to bet real money on.
\end{tryit}

% ══════════════════════════════════════════════════════════════════════════
\appendix
\chapter{Quick reference}
% ══════════════════════════════════════════════════════════════════════════

\section{Seed sizes (chained protocol 0.3.0)}
One mandatory text-only stage~0, then one 32-bit point per later stage
($N = \text{entropy bits}/32 = \text{words}/3$):

\begin{center}\small
\begin{tabular}{r r r r l}
\toprule
\textbf{Words} & \textbf{Entropy} & \textbf{Point stages $N$} & \textbf{Total (incl.\ 0)} & \textbf{Tier}\\
\midrule
3 & 32 & 1 & 2 & sub-standard\\
6 & 64 & 2 & 3 & sub-standard\\
9 & 96 & 3 & 4 & sub-standard\\
12 & 128 & 4 & 5 & standard (default)\\
15 & 160 & 5 & 6 & standard\\
18 & 192 & 6 & 7 & standard\\
21 & 224 & 7 & 8 & standard\\
24 & 256 & 8 & 9 & standard\\
\bottomrule
\end{tabular}
\end{center}
\noindent Hard cap: 256 bits / 24 words. Past that, add Argon2 iterations, not
stages. Default recommendation: 12 words.

\section{Glossary}
\begin{description}[leftmargin=0pt,style=nextline,nosep]
\item[TKBA] Tacit Knowledge-Based Authentication — the system's theoretical core.
\item[Tacit knowledge] Knowing-how you can deploy but not verbalize; the basis of
coercion-resistance.
\item[Burning Ship] The fractal whose structure is used as the shared, unstored
codebook.
\item[Bisection bijection] The lossless, reversible mapping from bits to a point
via recursive rectangle splits.
\item[I4F60] The 64-bit fixed-point format (60 fractional bits) that keeps the
math float-free and bit-exact.
\item[Argon2] The memory-hard function that gates each stage; the ``wall.''
\item[Stage 0] The mandatory text-only label that seeds the chain; salt or pepper.
\item[$N$] The Argon2 iteration count — the durable parameter you remember.
\item[$(o,p,q)$] The three numbers that perturb the fractal into your private one.
\item[TLP / RSW] Time-lock puzzle (repeated modular squaring); used to accelerate.
\item[jade-clock] The Lightning marketplace for outsourcing TLP solving.
\item[phoenix-scroll] The inheritance protocol (dead-man's switch + fallbacks).
\item[Graduation gate] The point at which training \emph{measures} that your
memory has consolidated and the setup is safe to rely on.
\end{description}

\section{The four properties, one more time}
\begin{enumerate}[leftmargin=1.6em,nosep]
\item Knowledge-based — in your head.
\item Individual — nobody else's.
\item Non-obscure — the attacker knows the method and that a stash exists.
\item Coercion-resistant — and still cannot rob it.
\end{enumerate}

\vfill
\begin{center}\small\itshape\color{graytext}
This is a pilot draft. Found an error, a confusing passage, or a better
analogy? That is the deliverable — please report it.
\end{center}

\end{document}
